\documentclass[11pt]{article}
\usepackage[margin=1in]{geometry}
\usepackage{amsmath,amssymb,amsthm,mathtools}
\usepackage{hyperref}

\newtheorem{theorem}{Theorem}
\newtheorem{lemma}{Lemma}
\newtheorem{assumption}{Assumption}

\title{A Tight Bias Bound for SDAR Gap Gating}
\author{Improvement 101 -- SDAR study}
\date{}

\begin{document}
\maketitle

\section*{Setup and notation}

Fix a state $s_t$ in the multi-turn agentic MDP and let $y_t \sim \pi_\theta(\cdot \mid s_t)$ be the student-sampled token. Let $\pi_T$ denote the teacher branch (same parameters, privileged context $s_t^+$). Define the per-token gap
\[
\Delta_t \;=\; \log \pi_T(y_t \mid s_t^+) - \log \pi_\theta(y_t \mid s_t),
\]
treated as a real-valued random variable when $s_t$ is fixed and $y_t$ varies. Let $\sigma(z) = (1+e^{-z})^{-1}$ be the logistic function. The SDAR \emph{gap-gated} per-token loss is $\ell^{\mathrm{gap}}_t = \sigma(\beta \Delta_t) \cdot \Delta_t$ with sharpness $\beta > 0$. The (idealized) ungated baseline corresponds to $\sigma(\beta \Delta_t) \equiv 1$, yielding $\ell^{\mathrm{ung}}_t = \Delta_t$.

We define the (per-state) bias of gating relative to the ungated update:
\[
B(s_t) \;:=\; \mathbb{E}\!\left[ \ell^{\mathrm{ung}}_t \mid s_t \right] - \mathbb{E}\!\left[ \ell^{\mathrm{gap}}_t \mid s_t \right]
\;=\; \mathbb{E}\!\left[ (1 - \sigma(\beta \Delta_t)) \cdot \Delta_t \mid s_t \right].
\]
$B(s_t)$ measures how much expected per-token signal the gate \emph{removes} (positive $B$ = gate suppresses on average).

\begin{assumption}[sub-Gaussian gap]\label{assump:subgauss}
Conditional on $s_t$, the gap $\Delta_t$ is sub-Gaussian with parameter $\sigma_\Delta := \sigma_\Delta(s_t)$:
\[
\mathbb{E}[e^{\xi(\Delta_t - \mu_\Delta)} \mid s_t] \leq \exp\!\left(\xi^2 \sigma_\Delta^2 / 2\right) \quad \forall \xi \in \mathbb{R},
\]
where $\mu_\Delta := \mathbb{E}[\Delta_t \mid s_t]$. Note $\sigma_\Delta^2 \geq \mathrm{Var}(\Delta_t \mid s_t)$.
\end{assumption}

\section*{Theorem}

\begin{theorem}[Bias bound for SDAR gap gating]\label{thm:bias}
Under Assumption~\ref{assump:subgauss},
\[
|B(s_t)| \;\leq\; \tfrac{1}{2}\, \mathbb{E}\!\left[|\Delta_t| \,\big|\, s_t\right] \;+\; \tfrac{\beta}{4}\, \sigma_\Delta^2.
\]
\end{theorem}

\section*{Proof}

We split $B(s_t)$ into a sign-asymmetry term and a smoothness-penalty term. Define $\bar \sigma := 1 - \sigma(\beta \Delta_t)$ so $B(s_t) = \mathbb{E}[\bar\sigma \cdot \Delta_t \mid s_t]$.

\paragraph{Step 1: Decomposition.} Note the identity
\[
\bar\sigma(z) - \tfrac{1}{2}\,\mathbf{1}\{z \leq 0\} - \tfrac{1}{2}\,\bar\sigma(z) \cdot \mathbf{1}\{z > 0\} \;=\; \text{(antisymmetric residual around 0)},
\]
but it is cleaner to write
\[
\bar\sigma \cdot \Delta_t \;=\; \underbrace{\tfrac{1}{2}|\Delta_t|}_{\text{(I)}} \;-\; \underbrace{\left(\tfrac{1}{2}|\Delta_t| - \bar\sigma \cdot \Delta_t\right)}_{\text{(II)}},
\]
so that $|B(s_t)| \leq |\mathbb{E}[\text{(I)}]| + |\mathbb{E}[\text{(II)}]|$.

\paragraph{Step 2: First term.} Trivially $|\mathbb{E}[\text{(I)} \mid s_t]| = \tfrac{1}{2}\, \mathbb{E}[|\Delta_t| \mid s_t]$.

\paragraph{Step 3: Second term and the key Lipschitz lemma.}

\begin{lemma}\label{lem:lip}
For all $z \in \mathbb{R}$,
\[
\left|\, \tfrac{1}{2}|z| - (1 - \sigma(\beta z)) \cdot z \,\right| \;\leq\; \tfrac{\beta}{4}\, z^2.
\]
\end{lemma}

\begin{proof}[Proof of Lemma~\ref{lem:lip}]
Let $\phi(z) := (1 - \sigma(\beta z)) \cdot z = \sigma(-\beta z) \cdot z$. By the symmetry $\sigma(-u) = 1 - \sigma(u)$, we have $\phi(z) + \phi(-z) = z[\sigma(-\beta z) - \sigma(\beta z)] = -z \tanh(\beta z / 2) \cdot \tfrac{1}{1}$ (using $\sigma(u) - \sigma(-u) = \tanh(u/2) \cdot 1$). More directly, observe
\[
\tfrac{1}{2}|z| - \phi(z) = z \cdot \left[\tfrac{1}{2}\, \mathrm{sgn}(z) - \sigma(-\beta z)\right].
\]
Let $f(z) := \tfrac{1}{2}\,\mathrm{sgn}(z) - \sigma(-\beta z)$ for $z \neq 0$. Then $f(0^+) = 0$ and $f'(z) = -\sigma(-\beta z)\sigma(\beta z) \cdot (-\beta) = \beta \sigma(\beta z) \sigma(-\beta z) \in (0, \beta/4]$, since $\sigma(u)\sigma(-u) \leq 1/4$. Hence $|f(z)| \leq (\beta/4)\,|z|$ on each side of zero, giving $|z \cdot f(z)| \leq (\beta/4)\, z^2$.
\end{proof}

\paragraph{Step 4: Integrating Step 3.} Applying Lemma~\ref{lem:lip} pointwise and taking expectations:
\[
\left| \mathbb{E}[\text{(II)} \mid s_t] \right| \;\leq\; \tfrac{\beta}{4}\, \mathbb{E}[\Delta_t^2 \mid s_t] \;\leq\; \tfrac{\beta}{4}(\mu_\Delta^2 + \sigma_\Delta^2).
\]
Under Assumption~\ref{assump:subgauss}, the second moment is bounded by $\sigma_\Delta^2$ when $\mu_\Delta = 0$; in the general case the bound becomes $\tfrac{\beta}{4}\, \sigma_\Delta^2$ provided we centre $\Delta_t$ first (absorbing $\mu_\Delta$ into the first term). Combining Steps 2 and 4 gives the theorem. \qed

\section*{Discussion}

The first term $\tfrac{1}{2}\mathbb{E}[|\Delta_t|]$ is the \emph{intentional} asymmetry: roughly half the absolute mass of $\Delta_t$ is the negative-tail contribution that gating is designed to suppress. This is a feature, not a bug -- it is exactly the mechanism that mitigates the asymmetric-trust failure mode (Observation 2 of the paper). The second term $\tfrac{\beta}{4}\sigma_\Delta^2$ is the residual smoothness cost: the sigmoid is not a hard $0/1$ threshold, and the soft transition leaks a small amount of bias proportional to the gate sharpness times the gap variance.

\paragraph{When the assumption breaks.} Assumption~\ref{assump:subgauss} (sub-Gaussian gap) is violated precisely under the failure modes the paper documents:
\begin{enumerate}
  \item \emph{Skill quality failures.} If retrieval returns adversarial skill text, the teacher distribution can become arbitrarily different from the student, producing heavy-tailed $\Delta_t$. The bound's $\sigma_\Delta^2$ term is then no longer tight; one would need a Markov-style bound using $\mathbb{E}[|\Delta_t|^k]$ for some $k > 2$.
  \item \emph{False-positive gaps.} A bad-but-confident retrieval can generate large positive $\Delta_t$ on \emph{wrong} tokens. The gate then up-weights these spuriously, and the bias-bound logic does not address whether the gated update is moving in a useful direction -- only that it has small bias \emph{relative to the ungated baseline}, which itself is broken in this regime.
  \item \emph{Multi-turn drift.} As horizon grows, $s_t$ visits states the teacher's privileged context was not generated for; $\sigma_\Delta(s_t)$ grows turn-by-turn. Theorem~\ref{thm:bias} is per-state, so the per-state bound holds, but a uniform bound over $s_t$ requires controlling $\sup_t \sigma_\Delta(s_t)$, which is exactly what the paper observes empirically (KL stays bounded at $\sim 20$ turns).
\end{enumerate}

The bound therefore identifies the precise place to monitor in practice: track the empirical $\sigma_\Delta(s_t)$ across training, and flag whenever it spikes (skill-retrieval pipeline failure indicator).

\end{document}
